% =========================================================
% Density of the Rational Numbers
% =========================================================

\subsection*{The Archimedean Property}
\label{sec:density-q}

\begin{exposition}
Two different ideas are often confused when one first studies the rational numbers.
The \emph{Archimedean property} says that the natural numbers grow beyond every
fixed rational. The \emph{density of $\mathbb{Q}$} says that between any two
distinct rationals there is another rational. The first controls large-scale
size; the second controls local refinement of intervals.
\end{exposition}

\begin{theorembox}{Theorem (Archimedean Property of $\mathbb{Q}$)}
\begin{theorem}[Archimedean property of $\mathbb{Q}$]
\label{thm:archimedean-q}
For every $x\in\mathbb{Q}$ there exists $n\in\mathbb{N}$ such that
\[
n>x.
\]
Equivalently, for every $x>0$ there exists $n\in\mathbb{N}$ such that
\[
\frac{1}{n}<x.
\]
\hyperref[prf:archimedean-q]{Go to proof.}
\end{theorem}
\end{theorembox}

\begin{remark*}[Standard quantified statement]
\[
\forall x\in\mathbb{Q}\;\exists n\in\mathbb{N}\;(n>x).
\]
\end{remark*}

\begin{remark*}[Interpretation]
The reciprocals of natural numbers become arbitrarily small. This theorem
rules out the existence of positive infinitesimals in $\mathbb{Q}$ and is the
foundation of all the approximation results in the chapter.
\end{remark*}

\begin{dependencies}
\begin{itemize}
  \item \hyperref[thm:q-is-ordered-field]{$\mathbb{Q}$ Is an Ordered Field}
  \item \hyperref[cor:n-positive]{Natural Numbers Are Positive}
\end{itemize}
\end{dependencies}

\begin{remark*}[Forward reference]
The statement that $\mathbb{Q}$ is dense in $\mathbb{R}$ is best treated as a
theorem about how the rationals sit inside the completed real line, not as an
internal theorem about $\mathbb{Q}$ alone. Its canonical home is therefore in
the reals chapter.
\end{remark*}
